\section{Conclusions}
\label{sec:conclusion}

In this project, we developed a configurable multi-stage Information Retrieval system for Task~1 of the CLEF~2026 CheckThat! Lab, focused on retrieving scientific sources for social media claims in English, French, and German. The system was designed as a modular pipeline separating parsing, text cleaning, text analysis, indexing, retrieval, and reranking, which allowed us to test different configurations systematically.

The results show that a carefully tuned lexical system remains a strong baseline for this task. Translating all inputs into English simplified the multilingual setting and enabled the use of a single English analyzer pipeline based on the Standard tokenizer, KStem, and an English stop list. BM25 was the strongest lexical ranking model among the alternatives tested, while fielded retrieval further improved performance by giving different weights to the title, abstract, and authors fields. In particular, the abstract field proved especially important, since social media claims usually paraphrase scientific findings rather than paper titles.

Our experiments also show that strict Boolean matching and standard lexical query expansion are not well suited to this task. AND constraints reduced recall sharply because claims and abstracts rarely share exact wording. Similarly, WordNet expansion and Pseudo-Relevance Feedback degraded performance due to ambiguity, vocabulary mismatch, and topic drift.

The best results were obtained by combining lexical and semantic evidence. Fusing BM25 scores with BGE-M3 dense similarity improved the average MRR@5 from 0.4719 for the best lexical system to 0.5290 for the hybrid system. This suggests that dense embeddings are most useful as a complementary signal to BM25, helping bridge the vocabulary and style gap between informal social media claims and formal scientific abstracts.

Our experiments with task-specific cross-encoder fine-tuning showed mixed results: the fine-tuned model improved French and German MRR@5 but degraded English, suggesting that the effect of fine-tuning is language-dependent and sensitive to training set size.